\documentclass[11pt]{article}
\usepackage{mathblatt}

\makeatletter
% Überschrift (Art, Gruppe) im Kasten der nächsten Aufgabe: nie allein am Seitenende (wie \pfab)
\let\gkopf\empty
\newcommand{\gueber}[1]{\gdef\gkopf{#1}}
\newcommand{\gueberdazu}[1]{\ifx\gkopf\empty\gdef\gkopf{#1}\else\expandafter\gdef\expandafter\gkopf\expandafter{\gkopf#1}\fi}
% eingerückt (B6): \gEin Einzug, \gSchrift eine Stufe kleiner; die Überschrift bleibt ohne Einzug
\newlength{\gEin}\setlength{\gEin}{0pt}
\let\gSchrift\relax
% Ankreuzzeile mit hängendem Einzug (lange Aussagen brechen unter dem Text um, nicht unter dem Kästchen)
\renewcommand{\pfkreuzzeile}[1]{\par\noindent\hangindent3.2em\hangafter1\hspace*{1.6em}$\square$\ #1\par}
\RenewDocumentEnvironment{pfaufg}{m m m}{%
  \begin{lrbox}{\mbaufgabebox}\begin{minipage}[t]{\linewidth}%
  \ifx\gkopf\empty\else\noindent\gkopf\par\global\let\gkopf\empty\fi
  \gSchrift\noindent\hspace*{\gEin}\mbpfrandmarke{#1}%
  \begin{minipage}[t]{\mbpfnr}\strut\textbf{#2}\end{minipage}%
  \begin{minipage}[t]{\dimexpr\linewidth-\gEin-\mbpfnr-\mbpfbe\relax}\raggedright\strut\ignorespaces}%
 {\end{minipage}%
  \begin{minipage}[t]{\mbpfbe}\raggedleft\small\color{mbgrau}\strut#3\end{minipage}%
  \end{minipage}\end{lrbox}%
  \setlength{\mbaufgabehoehe}{\dimexpr\ht\mbaufgabebox+\dp\mbaufgabebox\relax}%
  \par\addvspace{6pt}\Needspace*{\mbaufgabehoehe}%
  \noindent\usebox{\mbaufgabebox}\par\addvspace{6pt}}
\makeatother

\begin{document}
\pfheftstil
\fancyfoot[C]{\parbox[t]{\textwidth}{\mbfhbox\par\vspace{1.5mm}{\footnotesize\color{mbgrau}{\scriptsize Z8W}\hfill\thepage}}}
\pfheftkopf{Grundwert}{P10}
\begin{pfaufg}{}{1.}{}
Was ist hier das Ganze? Unterstreiche es. Rechne nicht.\par\smallskip\noindent a) In der Klasse sind $28$ Kinder. $7$ davon haben einen Hund.\par\noindent b) $45$ Gäste essen vegetarisch. Eingeladen sind $150$ Gäste.\par\noindent c) Ein Pullover kostete $60\,€$. Jetzt kostet er $45\,€$.
\end{pfaufg}
\par\addvspace{4pt}\begin{center}\begin{minipage}{0.5\linewidth}{\uebersichtskasten{\centering $G = W \cdot  100 : p$}}\end{minipage}\end{center}\par
\begin{pfaufg}{}{2.}{}
Schreibe zu jedem Satz auf, was W und was p ist. Schreibe dann die Rechnung für G mit den Zahlen hin. Rechne nicht aus.\par\smallskip\noindent a) 30 \% der Klasse sind 9 Kinder.\par\vspace{6mm}\noindent{\color{black!45}\rule{\linewidth}{0.4pt}}\par\smallskip\noindent b) 18 € Rabatt sind 15 \% des alten Preises.\par\vspace{6mm}\noindent{\color{black!45}\rule{\linewidth}{0.4pt}}
\end{pfaufg}
\par\addvspace{6pt}\gueber{\vspace{10pt}{\large\bfseries Vom Teil zum Ganzen}\par\vspace{2pt}}
\begin{pfaufg}{}{3.}{}
Der graue Abschnitt ist $20\,\%$ des Ganzen. Trage den Abschnitt so oft ab, bis du bei $100\,\%$ bist. Markiere das Ganze.
\par\smallskip\noindent \streifen[0]{20}{$0\,\%$}{}\par
\par\noindent \leerfeld[Abschnitte] \par
\end{pfaufg}
\begin{pfaufg}{}{4.}{}
Ein Händler hat $24$ kg Äpfel verkauft. Das sind $10\,\%$ seiner Äpfel. Wie viel kg Äpfel hatte er?
\par\noindent \leerfeld[kg] \par
\fusshilfe{\mbox{4:~$240$ kg}}
\end{pfaufg}
\begin{pfaufg}{}{5.}{}
$7\,\%$ eines Betrags sind $56\,€$. Wie viel ist das Ganze?
\par\smallskip\noindent \begin{tikzpicture}[baseline=(T.north)]\node[inner sep=0,outer sep=1.5pt] (T) {\renewcommand{\arraystretch}{1.35}\begin{tabular}{|r|r|}\hline \textbf{Prozent} & \textbf{€} \\ \hline $7\,\%$ & $56$\,€ \\ \hline $1\,\%$ & \leerzelle \\ \hline $100\,\%$ & \leerzelle \\ \hline\end{tabular}};\draw[->,mbgrau,line width=0.5pt] ($(T.north west)!0.3750!(T.south west)$) to[out=180,in=180,looseness=1.8] ($(T.north west)!0.6250!(T.south west)$);\draw[->,mbgrau,line width=0.5pt] ($(T.north east)!0.3750!(T.south east)$) to[out=0,in=0,looseness=1.8] ($(T.north east)!0.6250!(T.south east)$);\draw[->,mbgrau,line width=0.5pt] ($(T.north west)!0.6250!(T.south west)$) to[out=180,in=180,looseness=1.8] ($(T.north west)!0.8750!(T.south west)$);\draw[->,mbgrau,line width=0.5pt] ($(T.north east)!0.6250!(T.south east)$) to[out=0,in=0,looseness=1.8] ($(T.north east)!0.8750!(T.south east)$);\end{tikzpicture}\par
\fusshilfe{\mbox{5:~$800\,€$}}
\end{pfaufg}
\begin{pfaufg}{}{6.}{}
$36\,\%$ einer Menge sind $81$ kg. Wie viel kg ist die ganze Menge?
\par\noindent \leerfeld[kg] \par
\fusshilfe{\mbox{6:~$225$ kg}}
\end{pfaufg}
\begin{pfaufg}{nach P10 ’23}{7.}{}
Mia gibt 25 \% ihres Gewinns an ihre Eltern ab; das sind 200 €. Gib an, wie hoch der ganze Gewinn war.
\rechenplatz{1}
\fusshilfe{\mbox{7:~800 €}}
\end{pfaufg}
\par\setlength{\gEin}{8mm}\let\gSchrift\small
\begin{pfaufg}{P10 ’25}{8.}{}
Ein T-Shirt ist um 6 € billiger geworden. Dieser Nachlass entspricht 20 \% des ursprünglichen Preises. Gib den ursprünglichen Preis an.
\rechenplatz{1}
\fusshilfe{\mbox{8:~30 €}}
\end{pfaufg}
\begin{pfaufg}{nach P10 ’19}{9.}{}
In einem Gefäß liegen 4 schwarze Kugeln. Sie sind $\tfrac{2}{3}$ aller Kugeln im Gefäß. Berechne, wie viele Kugeln insgesamt im Gefäß liegen.
\rechenplatz{2}
\fusshilfe{\mbox{9:~G = 4 : $\tfrac{2}{3}$ = 6 Kugeln ($\tfrac{1}{3}$ sind 2 Kugeln)}}
\end{pfaufg}
\par\setlength{\gEin}{0pt}\let\gSchrift\relax
\par\vfill\noindent{\footnotesize\color{black!70}Grundwert in der P10: 3× seit 2014 (EBR und FOR). Nicht geprüft: vermehrter und verminderter Grundwert.\par\noindent Vorher: Prozent von einem Betrag · Weiter: Erhöhung und Senkung in Prozent\par}
\end{document}